\chapter{Integración con NOVA y trabajo siguiente}\label{cap:nova}

\section{Cómo encajaría en la nave}

\begin{cuidado}
Todo lo de esta sección está diseñado, pero no construido. El cambio en la nave va en un pull request aparte y lo
integra Nicolás Moreno.
\end{cuidado}

NOVA ya levanta sus modelos como servicios locales, que la propia nave abre y vigila: lenguaje en el puerto 8081, voz
en el 8082, oído en el 8083, búsqueda en el 8084 y visión en el 8792. El rango esperado y el pronóstico correrían como
un servicio Python más, en el 8793, que lee la historia de signos vitales de cada tripulante y devuelve el valor
esperado, el rango, el estado de la regla de 2 de 3 y el pronóstico a 7 días (\cref{fig:nave}).

\begin{figure}[htbp]
\centering
\begin{tikzpicture}[
  caja/.style={draw=#1, line width=0.8pt, rounded corners=2pt, fill=white, align=center, font=\footnotesize,
    text width=3.2cm, minimum height=1.3cm, inner sep=3pt},
  caja/.default=tinta3,
  flecha/.style={-{Stealth[length=2.2mm]}, line width=0.8pt, color=tinta2},
  nota/.style={font=\scriptsize\sffamily\color{tinta2}, align=center}]
  \node[caja] (est) at (0,2.4) {\textbf{Estación médica}\\ y teléfono de cada\\ tripulante};
  \node[caja=azulnoche] (nucleo) at (4.8,2.4) {\textbf{La nave}\\ F~Prime y enlace:\\ umbrales fijos y alertas};
  \node[caja=serieA, dashed] (serv) at (4.8,-0.6) {\textbf{Servicio :8793}\\ rango esperado\\ y pronóstico};
  \node[caja=tinta3] (ask) at (9.6,-0.6) {\textbf{Ask NOVA}\\ el modelo de lenguaje\\ explica, no decide};
  \node[caja=tinta3] (pant) at (9.6,2.4) {\textbf{Pantallas}\\ Vitals y NOVA Station};
  \node[caja=serieB] (tierra) at (0,-0.6) {\textbf{Tierra}\\ recibe el reporte\\ y firma paquetes};
  \draw[flecha] (est) -- node[nota, above] {signos} (nucleo);
  \draw[flecha] (nucleo) -- (pant);
  \draw[flecha] (nucleo.south west) -- node[nota, sloped, above] {reporte nocturno} (tierra.north east);
  \draw[flecha, dashed] (nucleo) -- node[nota, right, align=left] {historia\\ diaria} (serv);
  \draw[flecha, dashed] (serv) -- node[nota, above] {JSON} (ask);
  \draw[flecha, dashed] (serv.north east) -- node[nota, sloped, below] {tendencia a 7 días} (pant.south west);
  \draw[flecha, dashed] (tierra) -- node[nota, below, align=center] {modelo\\ base,\\ firmado} (serv);
  \node[nota, anchor=north] at (4.8,-1.4) {el corrector se reajusta\\ a bordo cada día};
\end{tikzpicture}
\caption[Diseño de la conexión con la nave]{Diseño de la conexión con la nave. Las líneas continuas existen hoy; las
punteadas están diseñadas, pero no construidas. Los umbrales fijos de emergencia no cambian.}
\label{fig:nave}
\end{figure}

Las decisiones que dan forma a ese diseño son estas (\cref{ap:decisiones}):
\begin{description}
\item[Dónde corre] Un servicio Python local en el puerto 8793, no un modelo exportado al enlace en Rust.
\item[Qué se adapta a bordo] El corrector son unos pocos coeficientes por variable y se puede reajustar a bordo cada
  día, sin esperar a la Tierra. El modelo base se mejora en la Tierra con los reportes nocturnos y vuelve a la nave
  como paquete firmado, el mismo mecanismo que hoy trae notas y respuestas corregidas.
\item[Qué no reemplaza] El rango adaptado no se propone, por ahora, como el rango \enquote{esperado en vuelo} de los
  umbrales de la nave. Los umbrales fijos de emergencia no se tocan, y el pronóstico no reemplaza el cono de la
  pantalla de signos vitales: va aparte, como una tendencia a 7 días.
\item[Quién explica] El servicio entrega números; el modelo de lenguaje de Ask NOVA los explica, sin generar
  predicciones propias ni diagnósticos. Esa conexión y qué variables llegan por la estación los define Nicolás Moreno.
\item[Qué queda fuera] La radiación queda diseñada como entrada del corrector (dosis diaria), sin construir; hay que
  definir con el equipo qué serie usar. El ánimo queda como variable opcional para más adelante, porque en NOVA solo
  existe si cada persona decide compartirlo.
\end{description}

\section{La corrida siguiente (v4)}\label{sec:v4}

Las respuestas del 14-11-2026 a las preguntas abiertas (\cref{ap:decisiones}) definieron una cuarta corrida, que va en
un pull request aparte:
\begin{enumerate}
\item \textbf{Solo pulso, presión sistólica y presión diastólica}, con los datos reales del reposo. Los escalones
  sintéticos se mantienen, declarados, porque son la única forma de medir la detección.
\item \textbf{LifeSnaps se mantiene y se mide cuánto aporta.} Sus ocho canales siguen entrando al pre-entrenamiento
  compartido; de las tres variables de la v4, solo el pulso está también en LifeSnaps.
\item \textbf{Todos los brazos en igualdad de condiciones}: misma tasa de aprendizaje y presupuesto, un tope de 60
  épocas o más para la cabeza, y 3 semillas.
\item \textbf{\enquote{Solo la cabeza} de verdad}, con el backbone en modo de evaluación, comparada con la versión
  actual (cabeza más BatchNorm).
\item \textbf{Pronóstico adaptado a cada persona}, con sus primeros $k = 7$, 14 y 28 días en cama.
\item \textbf{Ancho del rango por persona}, comparado con el poblacional.
\item \textbf{LoRA}, solo si el ajuste completo sigue ganándole al modelo desde cero con las 3 semillas.
\end{enumerate}

La auditoría deja además tareas de código para la misma corrida: aserciones de que prueba y entrenamiento no se cruzan,
la versión de Chronos fija y el hash del código en \texttt{metricas.json}. Para no repetir el problema de las muchas
comparaciones, conviene fijar de antemano los pares que responden cada pregunta. La demo del hackatón usa la presión y
el pulso, y los modelos de Hugging Face se harán públicos al final del hackatón, con sus fichas corregidas.
